%! AP Statistics — Unit 7, Lesson 7.1 — Guided Notes
\documentclass[10pt]{article}
\usepackage{apstats-article}
\usepackage{apstats-boxes}

\begin{document}

\pageheader{Unit 7, Lesson 7.1}{Guided Notes}
\namedateperiod

\begin{objectivebox}
By the end of this lesson, I will be able to:
\begin{itemize}
    \item[$\square$] \writeline
    \item[$\square$] \writeline
    \item[$\square$] \writeline
\end{itemize}
\end{objectivebox}

\begin{vocabbox}
\termblanklong{Type I error}
\termblanklong{Type II error}
\termblanklong{Significance level ($\alpha$)}
\termblanklong{$t$-distribution}
\end{vocabbox}

\begin{hookbox}
\textbf{Discussion:} A drug trial finds a statistically significant result. A follow-up
study contradicts it. How is this possible?

\writelines{3}
\end{hookbox}

\begin{notesbox}{Errors in Statistical Inference}
An inference procedure can produce two types of incorrect decisions:

\renewcommand{\arraystretch}{2}
\begin{tabularx}{\linewidth}{lXX}
    & \textbf{$H_0$ is True} & \textbf{$H_0$ is False} \\
    \hline
    \textbf{Reject $H_0$} & \blank{3cm} & Correct decision \\
    \textbf{Fail to reject $H_0$} & Correct decision & \blank{3cm} \\
\end{tabularx}

\vspace{0.2in}
\textbf{Type I error:} rejecting a \blank{3cm} null hypothesis. \\[4pt]
The probability of a Type I error $=$ \blank{2cm} (the significance level).

\vspace{0.15in}
\textbf{Type II error:} failing to reject a \blank{3cm} null hypothesis. \\[4pt]
The probability of a Type II error $= \beta$ (beta). \textbf{Power} $= 1 - \beta$.
\end{notesbox}

\begin{notesbox}{The Trade-off Between Error Types}
\textbf{Key Principle:} For a fixed sample size, decreasing the probability of one type
of error \blank{4cm} the probability of the other.

\vspace{0.1in}
If we lower $\alpha$ from 0.05 to 0.01:
\begin{itemize}
    \item Type I error rate \blank{2cm} (decreases / increases)
    \item Type II error rate \blank{2cm} (decreases / increases)
    \item Power \blank{2cm} (decreases / increases)
\end{itemize}

\vspace{0.1in}
To reduce \emph{both} error rates simultaneously, we must \blank{5cm}.
\end{notesbox}

\begin{notesbox}{Connecting to Unit 7 --- Why the $t$-Distribution?}
In Unit 6, inference for proportions used the \blank{2cm} distribution because
$\sigma_{\hat p} = \sqrt{p(1-p)/n}$ could be estimated from the hypothesized $p$.

For means, $\sigma$ is \blank{4cm}, so we substitute the sample standard deviation
$s$. This introduces extra variability, producing the \blank{3cm}-distribution.

\vspace{0.1in}
The $t$-statistic: \quad $t = \dfrac{\bar x - \mu_0}{s/\sqrt{n}}$ \quad follows a
$t$-distribution with $df = $ \blank{2cm} degrees of freedom.

\vspace{0.1in}
The $t$-distribution has \blank{4cm} tails than the standard normal (more / fewer).
As $df$ increases, the $t$-distribution approaches the \blank{3cm} distribution.
\end{notesbox}

\begin{practicebox}
\textbf{Guided Practice:} A hospital tests whether a new pain medication reduces
average recovery time below 14 days ($H_0{:}\ \mu = 14$, $H_a{:}\ \mu < 14$).

\begin{enumerate}
  \item Describe a Type I error in context.
        \writelines{2}

  \item Describe a Type II error in context.
        \writelines{2}

  \item Which error is more costly for patients? Explain.
        \writelines{2}

  \item If $\alpha$ is lowered from 0.05 to 0.01, describe the effect on both
        error types.
        \writelines{2}
\end{enumerate}
\end{practicebox}

\end{document}
